\section{Penurunan Penalti Batasan Portofolio}

\subsection{Suku Penalti Kuadratik}
Untuk memastikan sistem memilih tepat $K$ aset, digunakan penalti kuadratik dengan pengali Lagrange $A$:
\begin{equation}
P(x) = A \left(\sum_{i=1}^N x_i - K \right)^2
\label{eq:pen_qubo_derivation}
\end{equation}
Ekspansi binomial dari fungsi tersebut adalah:
\begin{equation}
P(x) = A \left( \left(\sum_{i=1}^N x_i\right)^2 - 2K\sum_{i=1}^N x_i + K^2 \right)
\end{equation}

\subsection{Ekspansi Suku Kuadratik Penjumlahan}
Suku $\left(\sum_{i=1}^N x_i\right)^2$ dijabarkan menjadi:
\begin{equation}
\left(\sum_{i=1}^N x_i\right)^2 = \sum_{i=1}^N x_i^2 + 2\sum_{i<j}^N x_i x_j
\end{equation}
Dengan memanfaatkan $x_i^2 = x_i$:
\begin{equation}
\left(\sum_{i=1}^N x_i\right)^2 = \sum_{i=1}^N x_i + 2\sum_{i<j}^N x_i x_j
\end{equation}

\subsection{Bentuk QUBO Penalti}
Substitusi kembali ke persamaan penalti $P(x)$:
\begin{equation}
\begin{aligned}
P(x) &= A \left( \sum_{i=1}^N x_i + 2\sum_{i<j}^N x_i x_j - 2K\sum_{i=1}^N x_i + K^2 \right) \\
&= A \left( (1-2K)\sum_{i=1}^N x_i + 2\sum_{i<j}^N x_i x_j + K^2 \right)
\end{aligned}
\end{equation}
Dalam format QUBO, suku konstanta $K^2$ dapat diabaikan dalam optimasi (offset energi).
\begin{equation}
H_{pen}(x) = \sum_{i=1}^N A(1-2K) x_i + \sum_{i < j} 2A x_i x_j
\end{equation}
Ini memungkinkan batasan $K$ diintegrasikan langsung ke dalam fungsi energi.
